%%====================================================================%%
%%                         INTRODUÇÃO                                  %%
%%====================================================================%%

\section{Introdução}
\label{sec:introducao}

Atualmente, sistemas de automação e monitoramento utilizam cada vez mais a comunicação em rede para permitir a troca de informações entre dispositivos físicos e aplicações de supervisão. Dentro desse contexto, uma das formas de realizar essa comunicação é através do protocolo HTTP associado a uma API REST, permitindo que diferentes dispositivos realizem operações de envio, consulta e alteração de informações através de uma rede.

Neste trabalho foi desenvolvido um sistema de comunicação entre uma placa ESP32 e uma interface de supervisão construída no Node-RED. O objetivo principal foi realizar a aquisição de variáveis ambientais utilizando um sensor BME280 e disponibilizar essas informações através de uma API REST. Além da aquisição dos sensores, também foi realizada a leitura de um botão físico e o controle remoto de um LED conectado ao ESP32.

O sensor BME280 permite realizar medições de temperatura, umidade relativa e pressão atmosférica. A partir da pressão também é possível realizar uma estimativa da altitude, utilizando como referência um valor conhecido de pressão ao nível do mar. Essas informações são adquiridas pelo ESP32 e posteriormente disponibilizadas para utilização pelo sistema.

A comunicação entre os diferentes elementos foi realizada através de uma rede Wi-Fi utilizando requisições HTTP. A API REST permite realizar a troca de informações entre o dispositivo e o Node-RED, fazendo com que os dados adquiridos pelo ESP32 possam ser apresentados em um dashboard e, no sentido contrário, que comandos realizados através da interface possam alterar o estado do LED.

O Node-RED foi utilizado para construção da interface principal de supervisão. No dashboard desenvolvido é possível acompanhar os valores provenientes do BME280, verificar o estado do botão, visualizar o estado do LED e realizar seu acionamento remotamente. Dessa maneira, foi possível implementar tanto o monitoramento quanto a atuação através da comunicação REST.

Outro requisito desenvolvido durante o trabalho foi a programação remota do ESP32 utilizando OTA (\textit{Over-the-Air}). Após uma primeira programação realizada através da conexão física, o recurso permite enviar novas versões do firmware pela rede Wi-Fi, reduzindo a necessidade de acesso físico ao dispositivo durante atualizações posteriores.

Além dos requisitos principais do trabalho, também foram desenvolvidas algumas funcionalidades extras, que são apresentadas a seguir.
Dessa forma, o desenvolvimento principal do trabalho está concentrado na integração entre ESP32, sensor BME280, API REST e Node-RED, juntamente com a implementação da atualização remota através de OTA. Os demais recursos desenvolvidos são apresentados separadamente como funcionalidades adicionais.



mas distribuídos e conectados.